\documentclass{article}
\usepackage{amsmath,amsthm}

\newtheorem{assumption}{Assumption}
\newtheorem{definition}{Definition}
\newtheorem{theorem}{Theorem}
\newtheorem{lemma}{Lemma}

\begin{document}

\begin{assumption}[Randomized treatment]\label{ass:randomized}
Let $(Y(0),Y(1),A,X)$ be observed under randomized treatment assignment with
$A \perp (Y(0),Y(1)) \mid X$ and $0 < \pi(X) < 1$.
\end{assumption}

\begin{definition}[Average treatment effect]\label{def:ate}
The target estimand is
\[
\tau = \mathbb{E}[Y(1)-Y(0)].
\]
\end{definition}

\begin{lemma}[Score unbiasedness]\label{lem:score-unbiased}
Under Assumption~\ref{ass:randomized}, the inverse-propensity score for
$\tau$ has mean zero after subtracting the estimand in Definition~\ref{def:ate}.
\end{lemma}
\begin{proof}
Condition on $X$ and use treatment randomization to replace observed outcomes
by the corresponding potential outcomes. The two conditional terms telescope.
\end{proof}

\begin{theorem}[Asymptotic normality of the ATE estimator]\label{thm:ate-clt}
Assume Assumption~\ref{ass:randomized}, finite second moments, and a consistent
propensity estimate. Then the inverse-propensity weighted estimator of $\tau$
is asymptotically normal at rate $\sqrt{n}$.
\end{theorem}
\begin{proof}
Decompose the estimator into the unbiased score in Lemma~\ref{lem:score-unbiased}
plus a negligible nuisance remainder. Apply a central limit theorem and Slutsky's theorem.
\end{proof}

\end{document}
